\begin{afterword}
\summaryhead

The author notes that several aspiring rationalists of the author's acquaintance grew up in cults or with a mentally ill parent, and that an Orthodox Jewish upbringing did the same work for the author: it broke an emotional trust in the sanity of everyone around. The essay argues that no one grows as a rationalist until that trust breaks, since otherwise the strangeness of ideas such as cryonics or many-worlds holds them back.

The author then recalls trusting Science at eighteen and believing that making a testable prediction was all that duty required, until it became clear that this had given no protection against a foolish theory. No procedure, the essay concludes, makes reasoning safe. Science demands too little, and Bayesian reasoning is young and hard to apply. Only after one's methods fail in a real disaster, and one keeps going anyway, does the real search begin. A postscript says the essay contains a deliberate flaw, which it does not identify.

\responsehead

The essay has the shape of a conversion narrative. Growth begins with a break from family and community. Hesitation about the group's beliefs is a symptom of social dependence. The teacher was once like the reader. The true path begins only after suffering. The last sentence makes failed parents and dead gods the qualification for the search. The essay anticipates how this sounds and answers with a link to the author's own post on cults, then relabels the reader's unease as ``nervously seeking reassurance.''

The central claim has no evidence behind it. The first sentence says the author does not ask and does not know. The sample is the acquaintances who volunteered dramatic childhoods. From that the essay states a law, ``Until this core emotional trust is broken, you don't start growing as a rationalist,'' and repeats it as ``I've never seen anyone.'' The ``strange ideas'' that broken trust is supposed to free the reader to accept are cryonics and many-worlds, the author's own positions. The essay never considers a reader who sees the logic and finds it wrong.

The warning is not applied to its author. ``There is no known procedure you can follow that makes your reasoning defensible,'' and ``the formal math is impossible to apply.'' A week earlier, many-worlds ``wins outright'' and is ``an obvious fact''; five days earlier, the math was simple enough to ``count the lines of code.'' The essay that says no method is safe goes on to tell the reader which discipline to adopt and which fields to study, ending with the author's own.

The history is overstated where it can be checked. Bayesian reasoning does not lack textbooks or elders; the author named Laplace and Jaynes as sources in 2007. (The essay's narrower point, that none of them teaches how to correct oneself, is fair.) The founding story leaves out the author's own 2007 account, in which Traditional Rationality was a help in climbing out of the hole, though it was not enough to get things right.

The postscript protects the essay against criticism. It claims an unnamed flaw, so any flaw a reader finds may confirm the author's foresight. The warning about clinging is aimed at flaws that are not real, so real ones still count, but the reader cannot tell whether a flaw they find is the one the author left in.

In short: a sermon that tells readers to break their trust in family, peers and science, gives no evidence that doing so helps, and exempts its author's own certainties.
\end{afterword}
